\subsection*{Experimental: Optional \emph{Where-Binding} for \texttt{within}}

You can append a \emph{where-clause} to impose \textbf{Where-Binding} constraints among different bounded variables defined in a \texttt{with-in} statement. This allows different bounded variables to be treated as aliases for each other. Below is an example:

\begin{lstlisting}[language=croqtile,caption={Experimental: Optional Where-Binding for within}]
with {m, n} in [M, N], {n_p, k} in [N_P, K]
where n_p <-> n {
  // matmul implements with m, n, K. 'n_p' always refers to the same value as 'n'.
}
\end{lstlisting}

Syntactically, the \texttt{\textless{}-\allowbreak{}\textgreater{}} operation establishes the \emph{where-binding} between two bounded variables. Within the \texttt{with-in} block, any reference to \texttt{n\_p} refers to \texttt{n}, and vice versa.
